1: % Reference entries, T (AN-11).  Law 1: os/iword/core/edit.c iw_ed_tab
2: % ([Tab] C1/02: to the next stop, every 0.5" relative to the left margin;
3: % WP recomputes the stop when the text before the tab changes, iw_ed_fixup;
4: % deleted without a question), spec/iword/wp5_format.h "tabs, indents and
5: % hard pages", face_a.c (Tab; Shift-, Ctrl- and Alt-Tab act as Tab).
6: % Shots: ../wordperfect51-decomp mint MINT09, MINT10, MINT21, MINT23.
7: % GATES: test-pubs-iword test-iword-codes test-iword-wp5-roundtrip test-iword-face test-flair-iword-block
8: \Gates{test-pubs-iword test-iword-codes test-iword-wp5-roundtrip test-iword-face test-flair-iword-block}
9: \begin{Feature}{Tab}
10: Tab moves the text that follows to the next tab stop.  Use it to indent the
11: first line of a paragraph and to line up short items in columns, such as
12: the entries of a form.
13: 
14: To tab,
15: \begin{WPSteps}
16: \item Press \K{Tab}.  The cursor moves to the next tab stop.
17: \item Type the text.
18: \end{WPSteps}
19: \Documents{wpkey}{Tab}
20: 
21: \BeginReplay{ref-tab}
22: For example, type \TypeRun{Name:}, press \PressKey{Tab}, type
23: \TypeRun{Peter Gibbons} and press \PressKey{Enter}; type \TypeRun{Dept:},
24: press \PressKey{Tab} twice\Keys{Tab} and type \TypeRun{Software}.  Press
25: \Press{Reveal Codes}{Alt-F3} to see the tab codes:
26: \WordScreen[rows=14]{ref-tab}{Tabs on the editing screen and in Reveal Codes}
27: \EndReplay
28: 
29: \Notes
30: \Topic{Code}
31: Each Tab puts a \Code{[Tab]} code in the document.  Delete it to remove the
32: tab; InitechWord deletes it without asking.
33: 
34: \Topic{Shift-Tab, Ctrl-Tab and Alt-Tab}
35: These act as Tab does.
36: 
37: \Topic{Tab Stops}
38: The tab stops are every half inch, measured from the left margin: under the
39: standard margins, every five columns of the editing screen.  You cannot set
40: other tab stops.  When you change the left margin the stops move with it.
41: Reveal Codes shows the stops on its ruler.
42: 
43: \Topic{Text Before a Tab}
44: If you add or delete text before a tab on the same line, the tab moves the
45: text after it to the next stop after the text, so tabbed columns stay lined
46: up as long as the text before them fits.
47: 
48: \SeeAlso{Indent; Margins, Left and Right; Reveal Codes}
49: \end{Feature}
50: 
51: \begin{Feature}{Template}
52: The function-key template lies above the function keys of your keyboard and
53: names the feature each key starts.  Each key has four lines on the template:
54: Ctrl and the key in red, Alt and the key in blue, Shift and the key in
55: green, the key alone in black.  The key's name is printed under its lines,
56: and the keys are in groups of four, as on the keyboard, with a legend of the
57: colors between the groups.
58: 
59: A line printed in grey names a WordPerfect feature that InitechWord does
60: not carry out.  Pressing that combination shows the same name in grey on the
61: status line and does nothing else; the next key you press clears it and acts
62: as usual (Appendix F lists these keys).  A line is blank where the
63: combination does nothing at all: Shift, Alt and Ctrl with F11 and F12.
64: 
65: The panel at the left of the template is a quick list of the keys you use
66: most that are not function keys: the hard page, hard return, tab, Cancel,
67: and the Home combinations for the top and end of the document, the start and
68: end of a line, and the place before the codes.
69: 
70: \Notes
71: \Topic{Ordering a Template}
72: Ask your Initech representative for part number 1062-0000-20.
73: 
74: \SeeAlso{Cursor Movement; Status Line}
75: \end{Feature}
